% Teoremas
\theoremstyle{plain}
\newtheorem{theorem}{Teorema}
\newtheorem{lemma}{Lema}
\providecommand*{\lemmaautorefname}{Lema}
\newtheorem{proposition}{Proposição}
\providecommand*{\propositionautorefname}{Proposição}
\newtheorem{corollary}{Corolário}
\providecommand*{\corollaryautorefname}{Corolário}
\newtheorem{conjecture}{Conjectura}
\providecommand*{\conjectureautorefname}{Conjectura}
\newtheorem{definition}{Definição}
\providecommand*{\definitionautorefname}{Definição}
\newtheorem{notation}{Notação}
\providecommand*{\notationautorefname}{Notação}
\newtheorem{remark}{Observação}
\providecommand*{\remarkautorefname}{Observação}
\newtheorem{example}{Exemplo}
\providecommand*{\exampleautorefname}{Exemplo}
\newtheorem{note}{Nota}
\providecommand*{\noteautorefname}{Nota}

% Listings
\lstset{
  language=C++,
  basicstyle=\ttfamily,
  keywordstyle=\color{blue},
  stringstyle=\color{verde},
  commentstyle=\color{red},
  extendedchars=true,
  showspaces=false,
  showstringspaces=false,
  numbers=left,
  numberstyle=\tiny,
  breaklines=true,
  backgroundcolor=\color{green!10},
  breakautoindent=true,
  fontadjust=false
}

% Cores e informações do PDF
\definecolor{blue}{RGB}{41,5,195}
\definecolor{verde}{rgb}{0,0.5,0}
\makeatletter
\hypersetup{
  pdftitle={\@title},
  pdfauthor={\@author},
  pdfcreator={LaTeX with unesp-abnTeX2},
  pdfkeywords={abnt}{latex}{abntex}{abntex2}{trabalho acadêmico},
  colorlinks=true,
  linkcolor=blue,
  citecolor=blue,
  filecolor=magenta,
  urlcolor=blue,
  bookmarksdepth=4
}
\makeatother

% Comandos globais
\everymath{\displaystyle}
\renewcommand{\sin}{\mathrm{sen}}
\renewcommand{\tan}{\mathrm{tg}}
\renewcommand{\csc}{\mathrm{cossec}}
\renewcommand{\cot}{\mathrm{cotg}}
\DeclareMathOperator{\posto}{\mathrm{posto}}
\DeclareMathOperator{\conv}{\mathrm{conv}}
\DeclareMathOperator{\diag}{\mathrm{diag}}
\DeclareMathOperator{\argmin}{\mathrm{arg}\min}
\DeclareMathOperator{\argmax}{\mathrm{arg}\max}

% Parágrafos
\setlength{\parindent}{2.0cm}
\setlength{\parskip}{0.2cm}

% Subfiguras
\newsubfloat{figure}
\providecommand*{\subfigureautorefname}{Subfigura}

% Atalhos
\newcommand{\FX}{\mathcal{F}(X)}
\newcommand{\F}{\ensuremath{\mathcal{F}}}
\newcommand{\R}{\ensuremath{\mathds{R}}}
\newcommand{\Aa}{[A]^{\alpha}}

% Citação em sobrescrito
\newcommand{\citep}[1]{\textsuperscript{\cite{#1}}}
